% GENERATED FILE. Edit canonical lessons, then run npm run generate:books.
% canonical-source-hash: fnv1a64:4e8cc563edfc0ff9
% canonical-lessons: JA-C425-tsuuyaku, JA-C425-hon-yaku, JA-C425-toiawaseru, JA-C425-tsuujiru, JA-C425-juuden, JA-R425-review-420-425, JA-R425-review-final-425

\chapter{Asking for Help in Another Language}
\label{ch:ja-asking-for-help-in-another-language}
\hlchaptermodality{425}

\begin{chapteropening}
\textbf{By the end of this chapter:} \emph{I can get help when language is a problem.}

Complete the last lesson of chapter 425: I can get help when language is a problem.
\end{chapteropening}

\section[\texorpdfstring{\ja{つうやく} (tsūyaku)}{つうやく (tsūyaku)}]{\ja{つうやく} (tsūyaku) — interpreting; an interpreter}

\label{lesson:JA-C425-tsuuyaku}



\begin{quote}
Before the new one: say the Japanese for to search all over, then the Japanese for fortunately.
\end{quote}

\subsection*{You'll want to know: \ja{つうやく}}
\textbf{\ja{つうやく}} — \emph{tsūyaku} — \textquotedblleft{}interpreting; an interpreter\textquotedblright{}.

\begin{quote}
\textbf{\ja{がいこくで} \ja{あにが} \ja{つうやくを} \ja{しました。}} — \emph{gaikoku de ani ga tsūyaku o shimashita} — \textquotedblleft{}Abroad, my brother did the interpreting.\textquotedblright{}
\end{quote}

\begin{cousinweb}
\textbf{\ja{つうやく}} (\emph{tsūyaku}) is Sino-Japanese (kango): two Chinese characters read with their on'yomi, readings borrowed from Chinese. They are \emph{tsū}, \textquotedblleft{}pass through\textquotedblright{} (Middle Chinese \emph{thuwng}) and \emph{yaku}, \textquotedblleft{}translate\textquotedblright{} (\emph{yek}) (KANJIDIC2; Middle Chinese in Baxter's transcription).
\end{cousinweb}

\begin{grammarlens}[title={What you've built}]
One more word for getting around.
\end{grammarlens}

\subsection*{Your turn}
\begin{itemize}
\raggedright
  \item \emph{Say it:} \emph{tsūyaku}
  \item \emph{Say it:} \emph{tsūyaku}, once more
  \item \emph{Say it:} \emph{gaikoku de ani ga tsūyaku o shimashita} — Abroad, my brother did the interpreting.
  \item \emph{Recall:} say the Japanese for to search all over, then the Japanese for fortunately, then say \emph{tsūyaku} again
\end{itemize}

\begin{tcolorbox}[breakable,colback=teal!4,colframe=teal!35!black,title={Before you move on}]
What does \ja{つうやく} mean? (\textquotedblleft{}Interpreting; an interpreter\textquotedblright{}.) Say it once more.
\end{tcolorbox}
\section[\texorpdfstring{\ja{ほんやく} (hon'yaku)}{ほんやく (hon'yaku)}]{\ja{ほんやく} (hon'yaku) — translation (of writing)}

\label{lesson:JA-C425-hon-yaku}



\begin{quote}
Before the new one: say the Japanese for fortunately, then the Japanese for interpreting.
\end{quote}

\subsection*{You'll want to know: \ja{ほんやく}}
\textbf{\ja{ほんやく}} — \emph{hon'yaku} — \textquotedblleft{}translation (of writing)\textquotedblright{}.

\begin{quote}
\textbf{\ja{この} \ja{てがみの} \ja{ほんやくを} \ja{おねがいします。}} — \emph{kono tegami no hon'yaku o onegai shimasu} — \textquotedblleft{}Please translate this letter for me.\textquotedblright{}
\end{quote}

\begin{cousinweb}
\textbf{\ja{ほんやく}} (\emph{hon'yaku}) is Sino-Japanese (kango): two Chinese characters read with their on'yomi, readings borrowed from Chinese. They are \emph{hon}, \textquotedblleft{}turn over\textquotedblright{} (Middle Chinese \emph{phjon}) and \emph{yaku}, \textquotedblleft{}translate\textquotedblright{} (\emph{yek}) (KANJIDIC2; Middle Chinese in Baxter's transcription).
\end{cousinweb}

\begin{grammarlens}[title={What you've built}]
One more word for getting around.
\end{grammarlens}

\subsection*{Your turn}
\begin{itemize}
\raggedright
  \item \emph{Say it:} \emph{hon'yaku}
  \item \emph{Say it:} \emph{hon'yaku}, once more
  \item \emph{Say it:} \emph{kono tegami no hon'yaku o onegai shimasu} — Please translate this letter for me.
  \item \emph{Recall:} say the Japanese for fortunately, then the Japanese for interpreting, then say \emph{hon'yaku} again
\end{itemize}

\begin{tcolorbox}[breakable,colback=teal!4,colframe=teal!35!black,title={Before you move on}]
What does \ja{ほんやく} mean? (\textquotedblleft{}Translation (of writing)\textquotedblright{}.) Say it once more.
\end{tcolorbox}
\section[\texorpdfstring{\ja{といあわせる} (toiawaseru)}{といあわせる (toiawaseru)}]{\ja{といあわせる} (toiawaseru) — to make an enquiry}

\label{lesson:JA-C425-toiawaseru}



\begin{quote}
Before the new one: say the Japanese for interpreting, then the Japanese for translation.
\end{quote}

\subsection*{You'll want to know: \ja{といあわせる}}
\textbf{\ja{といあわせる}} — \emph{toiawaseru} — \textquotedblleft{}to make an enquiry\textquotedblright{}.

\begin{quote}
\textbf{\ja{でんわで} \ja{えきに} \ja{といあわせました。}} — \emph{denwa de eki ni toiawasemashita} — \textquotedblleft{}I phoned the station to ask.\textquotedblright{}
\end{quote}

Polite \textbf{\ja{といあわせます}}, \emph{toiawasemasu}; an ichidan verb. Past (polite): \textbf{\ja{ちちが} \ja{りょかんに} \ja{といあわせました。}}, \emph{chichi ga ryokan ni toiawasemashita}.

\begin{cousinweb}
\textbf{\ja{といあわせる}} (\emph{toiawaseru}) is a native Japanese word (wago), inherited rather than borrowed from Chinese, built from \emph{toi}, \textquotedblleft{}asking\textquotedblright{} (from \emph{tou}) and \emph{awaseru}, \textquotedblleft{}to bring together\textquotedblright{}; its kanji were attached later as kun'yomi (JMdict).
\end{cousinweb}

\begin{grammarlens}[title={What you've built}]
One more word for getting around.
\end{grammarlens}

\subsection*{Your turn}
\begin{itemize}
\raggedright
  \item \emph{Say it:} \emph{toiawaseru}
  \item \emph{Say it:} \emph{toiawaseru}, once more
  \item \emph{Say it:} \emph{denwa de eki ni toiawasemashita} — I phoned the station to ask.
  \item \emph{Recall:} say the Japanese for interpreting, then the Japanese for translation, then say \emph{toiawaseru} again
\end{itemize}

\begin{tcolorbox}[breakable,colback=teal!4,colframe=teal!35!black,title={Before you move on}]
What does \ja{といあわせる} mean? (\textquotedblleft{}To make an enquiry\textquotedblright{}.) Say it once more.
\end{tcolorbox}
\section[\texorpdfstring{\ja{つうじる} (tsūjiru)}{つうじる (tsūjiru)}]{\ja{つうじる} (tsūjiru) — to be understood; to get through}

\label{lesson:JA-C425-tsuujiru}



\begin{quote}
Before the new one: say the Japanese for translation, then the Japanese for to make an enquiry.
\end{quote}

\subsection*{You'll want to know: \ja{つうじる}}
\textbf{\ja{つうじる}} — \emph{tsūjiru} — \textquotedblleft{}to be understood; to get through\textquotedblright{}.

\begin{quote}
\textbf{\ja{日本語}\ja{が} \ja{つうじて、あんしんしました。}} — \emph{nihongo ga tsūjite, anshin shimashita} — \textquotedblleft{}I was relieved that my Japanese was understood.\textquotedblright{}
\end{quote}

Polite \textbf{\ja{つうじます}}, \emph{tsūjimasu}; an ichidan verb. Past (polite): \textbf{\ja{あにの} \ja{日本語}\ja{が} \ja{つうじました。}}, \emph{ani no nihongo ga tsūjimashita}.

\begin{cousinweb}
\textbf{\ja{つうじる}} (\emph{tsūjiru}) is built from both layers of the vocabulary: the Sino-Japanese \emph{tsū}, \textquotedblleft{}pass through\textquotedblright{} (Middle Chinese \emph{thuwng}), an on'yomi borrowed from Chinese, made into a verb with the native ending \emph{-jiru}, a form of \emph{suru}, \textquotedblleft{}to do\textquotedblright{}, inherited rather than borrowed (JMdict; KANJIDIC2; Middle Chinese in Baxter's transcription).
\end{cousinweb}

\begin{grammarlens}[title={What you've built}]
One more word for getting around.
\end{grammarlens}

\subsection*{Your turn}
\begin{itemize}
\raggedright
  \item \emph{Say it:} \emph{tsūjiru}
  \item \emph{Say it:} \emph{tsūjiru}, once more
  \item \emph{Say it:} \emph{nihongo ga tsūjite, anshin shimashita} — I was relieved that my Japanese was understood.
  \item \emph{Recall:} say the Japanese for translation, then the Japanese for to make an enquiry, then say \emph{tsūjiru} again
\end{itemize}

\begin{tcolorbox}[breakable,colback=teal!4,colframe=teal!35!black,title={Before you move on}]
What does \ja{つうじる} mean? (\textquotedblleft{}To be understood; to get through\textquotedblright{}.) Say it once more.
\end{tcolorbox}
\section[\texorpdfstring{\ja{じゅうでん} (jūden)}{じゅうでん (jūden)}]{\ja{じゅうでん} (jūden) — charging a battery; recharging}

\label{lesson:JA-C425-juuden}



\begin{quote}
Before the new one: say the Japanese for to make an enquiry, then the Japanese for to be understood.
\end{quote}

\subsection*{You'll want to know: \ja{じゅうでん}}
\textbf{\ja{じゅうでん}} — \emph{jūden} — \textquotedblleft{}charging a battery; recharging\textquotedblright{}.

\begin{quote}
\textbf{\ja{えきで} \ja{でんわを} \ja{じゅうでんしました。}} — \emph{eki de denwa o jūden shimashita} — \textquotedblleft{}I charged my phone at the station.\textquotedblright{}
\end{quote}

\begin{cousinweb}
\textbf{\ja{じゅうでん}} (\emph{jūden}) is Sino-Japanese (kango): two Chinese characters read with their on'yomi, readings borrowed from Chinese. They are \emph{jū}, \textquotedblleft{}fill\textquotedblright{} (Middle Chinese \emph{tsyhuwng}) and \emph{den}, \textquotedblleft{}electricity\textquotedblright{} (\emph{denH}) (KANJIDIC2; Middle Chinese in Baxter's transcription).
\end{cousinweb}

\begin{grammarlens}[title={What you've built}]
One more word for getting around.
\end{grammarlens}

\subsection*{Your turn}
\begin{itemize}
\raggedright
  \item \emph{Say it:} \emph{jūden}
  \item \emph{Say it:} \emph{jūden}, once more
  \item \emph{Say it:} \emph{eki de denwa o jūden shimashita} — I charged my phone at the station.
  \item \emph{Recall:} say the Japanese for to make an enquiry, then the Japanese for to be understood, then say \emph{jūden} again
\end{itemize}

\begin{tcolorbox}[breakable,colback=teal!4,colframe=teal!35!black,title={Before you move on}]
What does \ja{じゅうでん} mean? (\textquotedblleft{}Charging a battery; recharging\textquotedblright{}.) Say it once more.
\end{tcolorbox}
\section[(dialogue)]{Looking back: from Receipts and Fees to Asking for Help in Another Language}

\label{lesson:JA-R425-review-420-425}



\begin{quote}
Say the Japanese for a handling fee and a receipt.
\end{quote}

\subsection*{Your turn}
Say these aloud:

\begin{itemize}
\raggedright
  \item a handling fee, a receipt, postage, to pay a bill, to refund
  \item to haggle, the list price, a price rise, a price cut, merchandise
  \item food in the Japanese style, food in the Western style, a set meal, sliced raw fish, grilled chicken on skewers
  \item a feast, paying the bill, splitting the bill, to treat someone, a full stomach
  \item a lost item, to get separated, to lose sight of, to search all over, fortunately
  \item interpreting, translation, to make an enquiry, to be understood, charging a battery
\end{itemize}

\begin{tcolorbox}[breakable,colback=teal!4,colframe=teal!35!black,title={Before you move on}]
A handling fee? (\textbf{\ja{てすうりょう}}, \emph{tesūryō}.) A receipt? (\textbf{\ja{りょうしゅうしょ}}, \emph{ryōshūsho}.) Postage? (\textbf{\ja{そうりょう}}, \emph{sōryō}.) To pay a bill? (\textbf{\ja{しはらう}}, \emph{shiharau}.) To refund? (\textbf{\ja{はらいもどす}}, \emph{haraimodosu}.) To haggle? (\textbf{\ja{ねぎる}}, \emph{negiru}.) The list price? (\textbf{\ja{ていか}}, \emph{teika}.) A price rise? (\textbf{\ja{ねあげ}}, \emph{neage}.) A price cut? (\textbf{\ja{ねさげ}}, \emph{nesage}.) Merchandise? (\textbf{\ja{しょうひん}}, \emph{shōhin}.) Food in the Japanese style? (\textbf{\ja{わしょく}}, \emph{washoku}.) Food in the Western style? (\textbf{\ja{ようしょく}}, \emph{yōshoku}.) A set meal? (\textbf{\ja{ていしょく}}, \emph{teishoku}.) Sliced raw fish? (\textbf{\ja{さしみ}}, \emph{sashimi}.) Grilled chicken on skewers? (\textbf{\ja{やきとり}}, \emph{yakitori}.) A feast? (\textbf{\ja{ごちそう}}, \emph{gochisō}.) Paying the bill? (\textbf{\ja{かんじょう}}, \emph{kanjō}.) Splitting the bill? (\textbf{\ja{わりかん}}, \emph{warikan}.) To treat someone? (\textbf{\ja{おごる}}, \emph{ogoru}.) A full stomach? (\textbf{\ja{まんぷく}}, \emph{manpuku}.) A lost item? (\textbf{\ja{おとしもの}}, \emph{otoshimono}.) To get separated? (\textbf{\ja{はぐれる}}, \emph{hagureru}.) To lose sight of? (\textbf{\ja{みうしなう}}, \emph{miushinau}.) To search all over? (\textbf{\ja{さがしまわる}}, \emph{sagashimawaru}.) Fortunately? (\textbf{\ja{さいわい}}, \emph{saiwai}.) Interpreting? (\textbf{\ja{つうやく}}, \emph{tsūyaku}.) Translation? (\textbf{\ja{ほんやく}}, \emph{hon'yaku}.) To make an enquiry? (\textbf{\ja{といあわせる}}, \emph{toiawaseru}.) To be understood? (\textbf{\ja{つうじる}}, \emph{tsūjiru}.) Charging a battery? (\textbf{\ja{じゅうでん}}, \emph{jūden}.)
\end{tcolorbox}
\section[(dialogue)]{Asking for Help in Another Language — once more}

\label{lesson:JA-R425-review-final-425}



\begin{quote}
Say the Japanese for interpreting and translation.
\end{quote}

\subsection*{Your turn}
\begin{itemize}
\raggedright
  \item \emph{Say it:} interpreting, translation, to make an enquiry, to be understood, charging a battery
\end{itemize}

\begin{tcolorbox}[breakable,colback=teal!4,colframe=teal!35!black,title={Before you move on}]
Interpreting? (\textbf{\ja{つうやく}}, \emph{tsūyaku}.) Translation? (\textbf{\ja{ほんやく}}, \emph{hon'yaku}.) To make an enquiry? (\textbf{\ja{といあわせる}}, \emph{toiawaseru}.) To be understood? (\textbf{\ja{つうじる}}, \emph{tsūjiru}.) Charging a battery? (\textbf{\ja{じゅうでん}}, \emph{jūden}.)
\end{tcolorbox}
